\documentclass[10pt,a4paper]{article}
\usepackage{../../styles/neet_style}

\pagestyle{fancy}
\fancyhf{}
\fancyhead[L]{\small \textbf{NEET UG 2027: Electrodynamics Master Theory}}
\fancyhead[R]{\small Satya Prakash}
\fancyfoot[L]{\small \textbf{NEET UG 2027 Prep Series}}
\fancyfoot[C]{\thepage}
\fancyfoot[R]{\small \textbf{Physics Chapter Notes}}

\begin{document}

% Title Header Banner
\begin{center}
  {\Huge \textbf{NEET UG 2027: Master Chapter Theory}}\\[6pt]
  {\Large \textbf{Electrodynamics, Magnetism, Electromagnetic Induction \& AC Circuits}}\\[4pt]
  {\small \textbf{Prepared, Compiled \& Authored by Satya Prakash} $\cdot$ Strictly Calibrated to Revised NMC Syllabus}
  \vspace{4pt}
  \hrule height 1.5pt
\end{center}

\section{Electrostatics, Coulomb's Law \& Continuous Charge Distributions}

\begin{conceptbox}[Foundations of Electric Charge \& Coulomb's Law in Vector Form]
\textbf{1. Fundamental Charge Properties:}
\begin{itemize}
  \item \textbf{Quantization of Charge:} Charge exists strictly as integer multiples of the elementary charge $e \approx 1.602 \times 10^{-19}\,\mathrm{C}$:
  \begin{equation}
    q = \pm n e \quad (n \in \mathbb{N})
  \end{equation}
  \item \textbf{Conservation of Charge:} The algebraic sum of fundamental charges within an isolated physical system remains invariant over time.
  \item \textbf{Invariance of Charge:} The magnitude of electric charge is independent of the reference frame speed ($q_{\text{rest}} = q_{\text{motion}}$), unlike relativistic mass.
\end{itemize}
\textbf{2. Coulomb's Inverse-Square Law in Vector Form:}
The electrostatic force exerted by point charge $q_1$ at position vector $\mathbf{r}_1$ on charge $q_2$ at $\mathbf{r}_2$:
\begin{equation}
  \mathbf{F}_{12} = \frac{1}{4\pi\varepsilon_0 \varepsilon_r} \frac{q_1 q_2}{|\mathbf{r}_2 - \mathbf{r}_1|^3} (\mathbf{r}_2 - \mathbf{r}_1) = \frac{1}{4\pi\varepsilon_0 \varepsilon_r} \frac{q_1 q_2}{r_{12}^2} \hat{\mathbf{r}}_{12}
\end{equation}
where $\varepsilon_0 = 8.854 \times 10^{-12}\,\mathrm{C}^2\cdot\mathrm{N}^{-1}\cdot\mathrm{m}^{-2}$ is the permittivity of free space, and $\frac{1}{4\pi\varepsilon_0} \approx 8.988 \times 10^9\,\mathrm{N}\cdot\mathrm{m}^2\cdot\mathrm{C}^{-2}$.
In a dielectric medium of relative permittivity (dielectric constant) $\kappa = \varepsilon_r > 1$:
\begin{equation}
  F_{\text{medium}} = \frac{F_{\text{vacuum}}}{\kappa}
\end{equation}
\textbf{3. Superposition Principle:}
The resultant force acting on any test charge $q_0$ due to a collection of discrete charges $q_1, q_2, \dots, q_n$ is the vector sum:
\begin{equation}
  \mathbf{F}_{\text{net}} = \sum_{i=1}^n \mathbf{F}_{i0} = \frac{q_0}{4\pi\varepsilon_0} \sum_{i=1}^n \frac{q_i}{|\mathbf{r}_0 - \mathbf{r}_i|^3} (\mathbf{r}_0 - \mathbf{r}_i)
\end{equation}
\end{conceptbox}

\begin{formulabox}[Electric Field \& Potential of Continuous Charge Geometries: Ring \& Disk]
\textbf{1. Electric Field Definition:}
$\mathbf{E} = \lim_{q_0 \to 0} \frac{\mathbf{F}}{q_0} = -\boldsymbol{\nabla} V$.
\textbf{2. Uniformly Charged Thin Circular Ring on its Symmetry Axis:}
Consider a ring of total charge $Q$ and radius $R$ positioned in the $y-z$ plane with symmetry axis along $x$.
\begin{itemize}
  \item By radial symmetry, perpendicular field components cancel: $dE_\perp = 0$.
  \item Axial field element: $dE_x = dE \cos\theta = \frac{1}{4\pi\varepsilon_0} \frac{dq}{x^2 + R^2} \frac{x}{\sqrt{x^2 + R^2}}$.
  \item Total integrated axial electric field:
  \begin{equation}
    \mathbf{E(x) = \frac{1}{4\pi\varepsilon_0} \frac{Q x}{(x^2 + R^2)^{3/2}} \hat{\mathbf{i}}}
  \end{equation}
  \item \textbf{At Ring Center ($x=0$):} $E(0) = 0$.
  \item \textbf{Far-Field Regime ($x \gg R$):} $E(x) \approx \frac{1}{4\pi\varepsilon_0}\frac{Q}{x^2}$ (acts as point charge).
  \item \textbf{Location of Maximum Electric Field:}
  \begin{equation}
    \frac{dE}{dx} = \frac{Q}{4\pi\varepsilon_0} \frac{(x^2 + R^2)^{3/2} - x \cdot \frac{3}{2}(x^2 + R^2)^{1/2}(2x)}{(x^2 + R^2)^3} = 0 \implies x^2 + R^2 - 3x^2 = 0 \implies \mathbf{x = \pm \frac{R}{\sqrt{2}}}
  \end{equation}
  \begin{equation}
    \mathbf{E_{\max} = \frac{1}{4\pi\varepsilon_0} \frac{2 Q}{3\sqrt{3} R^2}}
  \end{equation}
  \item \textbf{Electrostatic Potential on Ring Axis:}
  \begin{equation}
    \mathbf{V(x) = \frac{1}{4\pi\varepsilon_0} \frac{Q}{\sqrt{x^2 + R^2}}} \implies V_{\text{center}} = \frac{Q}{4\pi\varepsilon_0 R}
  \end{equation}
\end{itemize}
\textbf{3. Uniformly Charged Flat Circular Disk on its Axis:}
Integrating concentric rings of radius $r$ and charge $dq = \sigma (2\pi r dr)$ from $r=0$ to $R$:
\begin{equation}
  \mathbf{E(x) = \frac{\sigma}{2\varepsilon_0} \left(1 - \frac{x}{\sqrt{x^2 + R^2}}\right) \hat{\mathbf{i}}}, \qquad \mathbf{V(x) = \frac{\sigma}{2\varepsilon_0} \left(\sqrt{x^2 + R^2} - x\right)}
\end{equation}
For $R \to \infty$ (infinite sheet), $E = \frac{\sigma}{2\varepsilon_0}$ (constant, independent of $x$).
\end{formulabox}

\begin{formulabox}[Electric Dipole Dynamics: Field, Torque \& Potential Energy]
An electric dipole consists of two equal and opposite point charges $+q$ and $-q$ separated by a displacement vector $2\mathbf{a}$ pointing from $-q$ to $+q$:
\begin{equation}
  \mathbf{p} = q(2\mathbf{a}) \quad [\mathrm{C}\cdot\mathrm{m}]
\end{equation}
\textbf{1. Electric Field of Dipole in Coordinate Space ($r \gg a$):}
\begin{itemize}
  \item \textbf{Axial Point (End-on Position, $\theta = 0^\circ$):}
  \begin{equation}
    \mathbf{E_{\text{axial}} = \frac{1}{4\pi\varepsilon_0} \frac{2\mathbf{p}}{r^3}} \quad (\text{Collinear and parallel to } \mathbf{p})
  \end{equation}
  \item \textbf{Equatorial Point (Broadside-on Position, $\theta = 90^\circ$):}
  \begin{equation}
    \mathbf{E_{\text{eq}} = -\frac{1}{4\pi\varepsilon_0} \frac{\mathbf{p}}{r^3}} \quad (\text{Antiparallel to } \mathbf{p})
  \end{equation}
  \item \textbf{General Coordinate Point $(r, \theta)$:}
  \begin{equation}
    \mathbf{|\mathbf{E}| = \frac{1}{4\pi\varepsilon_0} \frac{p}{r^3} \sqrt{1 + 3\cos^2\theta}}, \qquad \tan\alpha = \frac{1}{2}\tan\theta
  \end{equation}
  where $\alpha$ is the angle between $\mathbf{E}$ and position vector $\mathbf{r}$.
  \item \textbf{Electrostatic Potential at $(r, \theta)$:}
  \begin{equation}
    \mathbf{V(r, \theta) = \frac{1}{4\pi\varepsilon_0} \frac{\mathbf{p}\cdot\hat{\mathbf{r}}}{r^2} = \frac{1}{4\pi\varepsilon_0} \frac{p\cos\theta}{r^2}}
  \end{equation}
  Notice: $V_{\text{eq}} = V(r, 90^\circ) = 0$ everywhere along the equatorial bisector plane.
\end{itemize}
\textbf{2. External Uniform Electric Field Interactions:}
\begin{itemize}
  \item Net Force: $\mathbf{F_{\text{net}} = q\mathbf{E} + (-q)\mathbf{E} = 0}$ (zero translational acceleration).
  \item Torque:
  \begin{equation}
    \mathbf{\boldsymbol{\tau} = \mathbf{p} \times \mathbf{E}} \implies |\boldsymbol{\tau}| = p E \sin\theta
  \end{equation}
  Maximum torque occurs at $\theta = 90^\circ$ ($\tau_{\max} = pE$).
  \item Potential Energy of Dipole in External Field:
  \begin{equation}
    \mathbf{U = -\mathbf{p}\cdot\mathbf{E} = -p E \cos\theta}
  \end{equation}
  \begin{itemize}
    \item \textbf{Stable Equilibrium ($\theta = 0^\circ$):} $\mathbf{p} \parallel \mathbf{E} \implies \boldsymbol{\tau} = 0$, $U_{\min} = -pE$.
    \item \textbf{Unstable Equilibrium ($\theta = 180^\circ$):} $\mathbf{p} \parallel -\mathbf{E} \implies \boldsymbol{\tau} = 0$, $U_{\max} = +pE$.
  \end{itemize}
  \item Work done by external agent in rotating dipole from $\theta_1$ to $\theta_2$:
  \begin{equation}
    \mathbf{W_{\text{ext}} = \Delta U = p E (\cos\theta_1 - \cos\theta_2)}
  \end{equation}
\end{itemize}
\end{formulabox}

\section{Gauss's Law, Electric Flux \& Boundary Value Applications}

\begin{conceptbox}[Gauss's Law Formulation \& Electric Flux]
\textbf{1. Electric Flux ($\Phi_E$):}
Surface integral of the electric field vector through an oriented surface element $d\mathbf{A} = \hat{\mathbf{n}} dA$:
\begin{equation}
  \Phi_E = \iint_S \mathbf{E}\cdot d\mathbf{A} \quad [\mathrm{N}\cdot\mathrm{m}^2\cdot\mathrm{C}^{-1} \text{ or } \mathrm{V}\cdot\mathrm{m}]
\end{equation}
\textbf{2. Gauss's Law (Maxwell's First Equation):}
The net outward electric flux through any closed Gaussian surface enclosing net charge $q_{\text{enclosed}}$:
\begin{equation}
  \mathbf{\oint_{\partial V} \mathbf{E}\cdot d\mathbf{A} = \frac{q_{\text{enclosed}}}{\varepsilon_0}}
\end{equation}
\textit{Key Physical Insights:}
\begin{itemize}
  \item Net flux is completely independent of the shape, size, or geometry of the Gaussian boundary.
  \item Charges residing outside the closed surface contribute zero net flux through the surface (entering flux equals exiting flux).
  \item Electric field $\mathbf{E}$ on the left-hand side is the \textbf{net resultant field} produced by all charges (both inside and outside).
\end{itemize}
\end{conceptbox}

\begin{formulabox}[Master Applications of Gauss's Law: Exact Field Derivations]
\textbf{1. Infinitely Long Line Charge of Uniform Density $\lambda$ ($\mathrm{C/m}$):}
Using a coaxial Gaussian cylinder of radius $r$ and length $L$:
\begin{equation}
  \oint \mathbf{E}\cdot d\mathbf{A} = E(2\pi r L) = \frac{\lambda L}{\varepsilon_0} \implies \mathbf{E(r) = \frac{\lambda}{2\pi\varepsilon_0 r}\hat{\mathbf{r}}}
\end{equation}
\textbf{2. Uniformly Charged Infinite Thin Flat Sheet of Density $\sigma$ ($\mathrm{C/m^2}$):}
Using a Gaussian pillbox of cross-sectional area $A$ piercing normal to the sheet:
\begin{equation}
  \oint \mathbf{E}\cdot d\mathbf{A} = 2 E A = \frac{\sigma A}{\varepsilon_0} \implies \mathbf{E = \frac{\sigma}{2\varepsilon_0}\hat{\mathbf{n}}} \quad (\text{Independent of distance})
\end{equation}
\textbf{3. Surface of a Charged Conductor:}
Inside the conductor, $E=0$. Flux exits only from the outer face:
\begin{equation}
  E A = \frac{\sigma A}{\varepsilon_0} \implies \mathbf{E_{\text{conductor}} = \frac{\sigma}{\varepsilon_0}\hat{\mathbf{n}}} \quad (\text{Twice that of a non-conducting sheet})
\end{equation}
\textbf{4. Thin Spherical Shell of Radius $R$ and Total Charge $Q$:}
\begin{itemize}
  \item \textbf{Interior Region ($r < R$):} Gaussian surface encloses $q_{\text{encl}} = 0 \implies \mathbf{E_{\text{in}} = 0}$.
  \item Potential is constant throughout interior: $\mathbf{V_{\text{in}} = \frac{1}{4\pi\varepsilon_0}\frac{Q}{R}}$.
  \item \textbf{Exterior Region ($r \ge R$):} Gaussian sphere encloses $Q \implies \mathbf{E_{\text{out}}(r) = \frac{1}{4\pi\varepsilon_0}\frac{Q}{r^2}\hat{\mathbf{r}}}$.
  \item Exterior potential: $\mathbf{V_{\text{out}}(r) = \frac{1}{4\pi\varepsilon_0}\frac{Q}{r}}$.
\end{itemize}
\textbf{5. Uniformly Charged Non-Conducting Solid Sphere of Radius $R$ and Total Charge $Q$:}
\begin{itemize}
  \item Volume charge density: $\rho = \frac{Q}{\frac{4}{3}\pi R^3}$.
  \item \textbf{Interior Region ($r \le R$):} $q_{\text{encl}} = \rho \left(\frac{4}{3}\pi r^3\right) = Q \frac{r^3}{R^3}$:
  \begin{equation}
    E(r) \cdot 4\pi r^2 = \frac{Q r^3}{\varepsilon_0 R^3} \implies \mathbf{E_{\text{in}}(r) = \frac{1}{4\pi\varepsilon_0}\frac{Q r}{R^3}\hat{\mathbf{r}} = \frac{\rho r}{3\varepsilon_0}\hat{\mathbf{r}}}
  \end{equation}
  Electric field increases linearly with radial distance $r$ from center.
  \item \textbf{Interior Potential:} $V_{\text{in}}(r) = -\int_\infty^R E_{\text{out}} dr - \int_R^r E_{\text{in}} dr$:
  \begin{equation}
    \mathbf{V_{\text{in}}(r) = \frac{Q}{4\pi\varepsilon_0}\left[\frac{3R^2 - r^2}{2R^3}\right]} \implies \mathbf{V_{\text{center}} = \frac{3}{2} V_{\text{surface}} = 1.5 V_{\text{surface}}}
  \end{equation}
  \item \textbf{Exterior Region ($r \ge R$):} Identical to point charge:
  \begin{equation}
    \mathbf{E_{\text{out}}(r) = \frac{1}{4\pi\varepsilon_0}\frac{Q}{r^2}\hat{\mathbf{r}}}, \qquad \mathbf{V_{\text{out}}(r) = \frac{1}{4\pi\varepsilon_0}\frac{Q}{r}}
  \end{equation}
\end{itemize}
\end{formulabox}

\begin{ncertbox}[NCERT Extraction: Electrostatic Shielding \& Cavity Boundary Properties]
Inside any hollow conducting cavity of arbitrary shape, the electric field is \textbf{identically zero} ($\mathbf{E} = 0$), provided no charge is enclosed within the cavity. The potential throughout the cavity remains constant and strictly equals the potential on the conducting surface ($V_{\text{cavity}} = V_{\text{surface}}$).
\begin{itemize}
  \item \textit{Faraday Cage Principle:} Any sensitive electrical apparatus enclosed in a hollow metallic enclosure is completely shielded from external electrical disturbances.
  \item \textit{Practical Safety:} During lightning storms, it is significantly safer to remain inside a closed metallic car or bus than under a tree or in the open, as lightning charges flow along the outer conducting body directly to the ground without penetrating the passenger interior.
\end{itemize}
\end{ncertbox}

\begin{trapbox}[Syllabus Boundary Alert: Van de Graaff Generator Deletion]
\textbf{NMC REVISION WARNING FOR NEET UG:}
Under the revised NEET UG syllabus issued by the National Medical Commission:
\begin{itemize}
  \item \textbf{STRICTLY DELETED:} The \textbf{Van de Graaff Generator} (high voltage production via belt charge transport, corona discharge, insulating column, spherical dome collection) has been \textbf{completely eliminated}.
  \item \textbf{EXAM ORIENTATION:} Candidates must master Gauss's law boundary applications (spherical shells, solid spheres, line charges, and conductor electrostatic shielding) without studying the obsolete Van de Graaff machine.
\end{itemize}
\end{trapbox}

\section{Capacitance, Dielectric Physics \& Energy Storage}

\begin{conceptbox}[Dielectric Polarization \& Capacitance Principles]
\textbf{1. Capacitance Definition:}
The capacity of a conductor to store electric charge per unit potential rise:
\begin{equation}
  C = \frac{Q}{V} \quad [\mathrm{Farad}, \mathrm{F} = \mathrm{C}\cdot\mathrm{V}^{-1}]
\end{equation}
For an isolated conducting sphere of radius $R$: $C = 4\pi\varepsilon_0 R$.
\textbf{2. Parallel Plate Capacitor in Vacuum:}
Plates of area $A$ separated by vacuum gap $d$ with uniform field $E = \frac{\sigma}{\varepsilon_0} = \frac{Q}{\varepsilon_0 A}$ and $V = Ed$:
\begin{equation}
  \mathbf{C_0 = \frac{\varepsilon_0 A}{d}}
\end{equation}
\textbf{3. Dielectric Material Polarization:}
When an insulating dielectric is inserted into an external electric field $\mathbf{E}_0$, bound molecular dipoles align, generating a macroscopic polarization vector $\mathbf{P} = \varepsilon_0 \chi_e \mathbf{E}$.
This induces surface bound charge density $\sigma_b = \mathbf{P}\cdot\hat{\mathbf{n}}$ producing an internal opposing depolarizing field $\mathbf{E}_{\text{induced}} = \frac{\sigma_b}{\varepsilon_0}$.
\begin{equation}
  \mathbf{E_{\text{net}} = \mathbf{E}_0 - \mathbf{E}_{\text{induced}} = \frac{\mathbf{E}_0}{\kappa}} \quad (\text{Field is attenuated by factor } \kappa)
\end{equation}
where $\kappa = 1 + \chi_e = \varepsilon_r$ is the dielectric constant.
\textbf{4. Bound Induced Surface Charge:}
\begin{equation}
  E_{\text{net}} = \frac{\sigma - \sigma_b}{\varepsilon_0} = \frac{\sigma}{\kappa\varepsilon_0} \implies \sigma_b = \sigma\left(1 - \frac{1}{\kappa}\right) \implies \mathbf{q_b = q\left(1 - \frac{1}{\kappa}\right)}
\end{equation}
\end{conceptbox}

\begin{formulabox}[Capacitor with Partial Dielectric Slab \& Network Combinations]
\textbf{1. Partial Dielectric Slab Insertion (Thickness $t < d$):}
Across the vacuum gap $(d-t)$, electric field is $E_0 = \frac{Q}{\varepsilon_0 A}$. Inside the slab of thickness $t$, field is $E = \frac{E_0}{\kappa}$.
Total potential difference between plates:
\begin{equation}
  V = E_0 (d - t) + E t = E_0 \left[(d - t) + \frac{t}{\kappa}\right] = \frac{Q}{\varepsilon_0 A} \left[d - t\left(1 - \frac{1}{\kappa}\right)\right]
\end{equation}
\begin{equation}
  \mathbf{C = \frac{Q}{V} = \frac{\varepsilon_0 A}{d - t\left(1 - \frac{1}{\kappa}\right)}}
\end{equation}
\textit{Special Cases:}
\begin{itemize}
  \item Complete dielectric filling ($t=d$): $C = \kappa \frac{\varepsilon_0 A}{d} = \kappa C_0$.
  \item Conducting slab inserted ($t < d$, $\kappa \to \infty$): $C = \frac{\varepsilon_0 A}{d - t}$.
\end{itemize}
\textbf{2. Grouping of Capacitors:}
\begin{itemize}
  \item \textbf{Series Combination (Charge $Q$ is identical):}
  \begin{equation}
    V = V_1 + V_2 + \dots \implies \mathbf{\frac{1}{C_{\text{eq}}} = \sum_{i=1}^n \frac{1}{C_i}} \quad \left(\text{For two: } C_{\text{eq}} = \frac{C_1 C_2}{C_1 + C_2}\right)
  \end{equation}
  \item \textbf{Parallel Combination (Potential difference $V$ is identical):}
  \begin{equation}
    Q = Q_1 + Q_2 + \dots \implies \mathbf{C_{\text{eq}} = \sum_{i=1}^n C_i}
  \end{equation}
\end{itemize}
\textbf{3. Electrostatic Energy Stored \& Field Energy Density:}
Work required to assemble charge $Q$ on plates against potential $V$:
\begin{equation}
  \mathbf{U = \int_0^Q V(q) dq = \int_0^Q \frac{q}{C} dq = \frac{Q^2}{2C} = \frac{1}{2} C V^2 = \frac{1}{2} Q V}
\end{equation}
\textbf{4. Energy Density in Electric Field:}
Dividing energy $U$ by the volume of the dielectric space ($A \cdot d$):
\begin{equation}
  u_E = \frac{U}{Ad} = \frac{\frac{1}{2}(\kappa \varepsilon_0 A/d)(Ed)^2}{Ad} \implies \mathbf{u_E = \frac{1}{2} \kappa \varepsilon_0 E^2}
\end{equation}
\textbf{5. Attractive Mechanical Force Between Plates:}
\begin{equation}
  \mathbf{F = -\frac{dU}{dx} = \frac{Q^2}{2\varepsilon_0 A} = \frac{1}{2} Q E_0 = \frac{1}{2} \varepsilon_0 A E_0^2}
\end{equation}
\end{formulabox}

\begin{usecasebox}[Dielectric Insertion: Battery Connected vs. Battery Disconnected]
\begin{center}
\begin{tabularx}{\linewidth}{|l|Y|Y|}
\hline
\textbf{Physical Parameter} & \textbf{Case A: Battery Disconnected ($Q = \text{const}$)} & \textbf{Case B: Battery Remains Connected ($V = \text{const}$)} \\
\hline
\textbf{Charge ($Q$)} & Unchanged: $Q = Q_0$ & Increases: $Q = \kappa Q_0$ \\
\hline
\textbf{Capacitance ($C$)} & Increases: $C = \kappa C_0$ & Increases: $C = \kappa C_0$ \\
\hline
\textbf{Potential Difference ($V$)} & Decreases: $V = V_0 / \kappa$ & Unchanged: $V = V_0$ \\
\hline
\textbf{Electric Field ($E$)} & Decreases: $E = E_0 / \kappa$ & Unchanged: $E = E_0$ \\
\hline
\textbf{Stored Energy ($U$)} & Decreases: $U = U_0 / \kappa$ & Increases: $U = \kappa U_0$ \\
\hline
\textbf{Battery Work Done} & Zero (battery isolated) & $W_{\text{battery}} = \Delta Q \cdot V_0 = (\kappa - 1) C_0 V_0^2$ \\
\hline
\end{tabularx}
\end{center}
\end{usecasebox}

\section{Current Electricity, Drift Dynamics \& DC Resistive Networks}

\begin{conceptbox}[Drift Velocity, Microscopic Ohm's Law \& Mobility]
\textbf{1. Microscopic Mechanism of Conduction in Metals:}
Under zero external field, conduction electrons execute chaotic thermal motion at speeds $v_{\text{th}} \sim 10^5\,\mathrm{m/s}$ with random directions ($\langle \mathbf{v}_{\text{th}} \rangle = 0$).
When an external electric field $\mathbf{E}$ is applied, each electron experiences an electrostatic force $\mathbf{F} = -e\mathbf{E}$, imparting constant acceleration $\mathbf{a} = -\frac{e\mathbf{E}}{m_e}$.
Electrons repeatedly scatter elastically from vibrating crystal lattice ions. The average time elapsed between consecutive scattering collisions is the \textbf{relaxation time} $\tau \sim 10^{-14}\,\mathrm{s}$.
\textbf{2. Drift Velocity ($\mathbf{v}_d$) Derivation:}
\begin{equation}
  \mathbf{v}_d = \langle \mathbf{v}_{\text{initial}} + \mathbf{a} t \rangle = 0 + \mathbf{a}\tau \implies \mathbf{\mathbf{v}_d = -\frac{e \mathbf{E} \tau}{m_e}}
\end{equation}
Typically, drift speed is remarkably slow: $v_d \sim 10^{-4}\,\mathrm{m/s} = 0.1\,\mathrm{mm/s}$.
\textbf{3. Current Density \& Microscopic Ohm's Law:}
Consider a conductor of cross-sectional area $A$ with free electron density $n$. In time $\Delta t$, charge passing through is $\Delta Q = n e A v_d \Delta t$.
\begin{equation}
  I = \frac{\Delta Q}{\Delta t} = \mathbf{n e A v_d} \implies \mathbf{J = \frac{I}{A} = n e v_d = \left(\frac{n e^2 \tau}{m_e}\right) E}
\end{equation}
\begin{equation}
  \mathbf{\mathbf{J} = \sigma \mathbf{E} = \frac{1}{\rho}\mathbf{E}} \quad (\text{Vector / Microscopic Ohm's Law})
\end{equation}
where electrical conductivity $\sigma$ and resistivity $\rho$:
\begin{equation}
  \mathbf{\sigma = \frac{n e^2 \tau}{m_e}}, \qquad \mathbf{\rho = \frac{m_e}{n e^2 \tau}}
\end{equation}
\textbf{4. Electron Mobility ($\mu$):}
Drift speed per unit electric field:
\begin{equation}
  \mathbf{\mu = \frac{v_d}{E} = \frac{e \tau}{m_e}} \quad [\mathrm{m}^2\cdot\mathrm{V}^{-1}\cdot\mathrm{s}^{-1}]
\end{equation}
\end{conceptbox}

\begin{formulabox}[Temperature Dependence, Cell Grouping \& Power Transfer Theorem]
\textbf{1. Temperature Dependence of Resistance:}
For metallic conductors, increasing temperature intensifies lattice vibrations, reducing relaxation time $\tau$:
\begin{equation}
  R(T) = R_0 [1 + \alpha (T - T_0)] \implies \mathbf{\alpha = \frac{R_T - R_0}{R_0 \Delta T}} \quad [^\circ\mathrm{C}^{-1} \text{ or } \mathrm{K}^{-1}]
\end{equation}
For metals, $\alpha > 0$ (positive). For semiconductors and electrolytes, charge carrier density surges exponentially with temperature ($n(T) \propto e^{-E_g / 2k_B T}$), overriding $\tau$ drops, giving $\mathbf{\alpha < 0}$ (negative).
\textbf{2. Terminal Potential Difference of a Real Cell ($\mathcal{E}, r$):}
\begin{itemize}
  \item \textbf{Discharging (Supplying current $I$ to load $R$):} $\mathbf{V = \mathcal{E} - I r} = \frac{\mathcal{E} R}{R + r}$.
  \item \textbf{Charging (External charger driving current into cell):} $\mathbf{V = \mathcal{E} + I r}$.
  \item \textbf{Open Circuit ($I=0$):} $V = \mathcal{E}$. Short circuit ($R=0$): $I_{\text{sc}} = \frac{\mathcal{E}}{r}, V = 0$.
\end{itemize}
\textbf{3. Mixed Grouping of Cells ($m$ rows of $n$ series cells each):}
Total cells $N = n \times m$. Total EMF $= n\mathcal{E}$; total internal resistance $= \frac{n r}{m}$.
Current through load $R$:
\begin{equation}
  I = \frac{n\mathcal{E}}{R + \frac{n r}{m}} = \frac{m n \mathcal{E}}{m R + n r}
\end{equation}
To maximize current, minimize the denominator $(mR + nr) = (\sqrt{mR} - \sqrt{nr})^2 + 2\sqrt{mnRr}$.
The minimum occurs when $\sqrt{mR} = \sqrt{nr}$:
\begin{equation}
  \mathbf{R = \frac{n r}{m}} \quad (\text{Load resistance equals total internal resistance})
\end{equation}
\textbf{4. Maximum Power Transfer Theorem:}
Power delivered to load: $P(R) = I^2 R = \frac{\mathcal{E}^2 R}{(R + r)^2}$. Setting $\frac{dP}{dR} = 0$:
\begin{equation}
  \frac{dP}{dR} = \mathcal{E}^2 \frac{(R + r)^2 - 2R(R + r)}{(R + r)^4} = 0 \implies R + r - 2R = 0 \implies \mathbf{R = r}
\end{equation}
\begin{equation}
  \mathbf{P_{\max} = \frac{\mathcal{E}^2}{4r}}
\end{equation}
\end{formulabox}

\begin{formulabox}[Kirchhoff's Laws \& Measurement Bridges]
\textbf{1. Kirchhoff's Laws:}
\begin{itemize}
  \item \textbf{Junction Rule (KCL):} At any circuit node: $\sum I_{\text{in}} = \sum I_{\text{out}}$ (Conservation of Charge).
  \item \textbf{Loop Rule (KVL):} Around any closed loop: $\sum \Delta V = \sum \mathcal{E} - \sum I R = 0$ (Conservation of Energy).
\end{itemize}
\textbf{2. Wheatstone Bridge Balance Derivation:}
A quadrilateral network of four resistors $P, Q, R, S$ with a galvanometer connected between junctions $B$ and $D$.
At balance, no current traverses the galvanometer ($I_g = 0 \implies V_B = V_D$).
Applying KVL to loop $ABDA$: $I_1 P - I_2 R = 0 \implies I_1 P = I_2 R$.
Applying KVL to loop $BCDB$: $I_1 Q - I_2 S = 0 \implies I_1 Q = I_2 S$.
Dividing the two equations yields the classical balance condition:
\begin{equation}
  \mathbf{\frac{P}{Q} = \frac{R}{S}}
\end{equation}
\textbf{3. Metre Bridge (Slide-Wire Bridge):}
A uniform wire of length $100\,\mathrm{cm}$ acts as ratio arms $P = \rho_w l$ and $Q = \rho_w (100 - l)$.
\begin{equation}
  \frac{R}{X} = \frac{l}{100 - l} \implies \mathbf{X = R \left(\frac{100 - l}{l}\right)}
\end{equation}
Maximum experimental precision occurs at the central null point ($l \approx 50\,\mathrm{cm}$).
\end{formulabox}

\begin{trapbox}[Syllabus Boundary Alert: Potentiometer \& Carbon Resistor Color Codes Deletion]
\textbf{NMC REVISION WARNING FOR NEET UG:}
Under the revised NEET UG syllabus issued by the National Medical Commission:
\begin{itemize}
  \item \textbf{STRICTLY DELETED:} The \textbf{Potentiometer} (principle of fall of potential $V \propto l$, comparison of EMFs of two cells $\mathcal{E}_1/\mathcal{E}_2 = l_1/l_2$, and determination of cell internal resistance $r = R(l_1/l_2 - 1)$) has been \textbf{completely deleted}.
  \item \textbf{STRICTLY DELETED:} \textbf{Color Coding of Carbon Resistors} (BBROYGBVGW mnemonic bands, tolerance multipliers) has been \textbf{completely deleted}.
  \item \textbf{RETAINED:} The \textbf{Metre Bridge} and \textbf{Wheatstone Bridge} are fully retained and form the exclusive focus for bridge measurements.
\end{itemize}
\end{trapbox}

\section{Magnetic Effects of Electric Current \& Moving Charges}

\begin{conceptbox}[Biot-Savart Law \& Magnetic Lorentz Force]
\textbf{1. Biot-Savart Law:}
The infinitesimal magnetic field $d\mathbf{B}$ generated by an elementary current length vector $I d\mathbf{l}$ at displacement $\mathbf{r}$:
\begin{equation}
  \mathbf{d\mathbf{B} = \frac{\mu_0}{4\pi} \frac{I (d\mathbf{l} \times \hat{\mathbf{r}})}{r^2} = \frac{\mu_0}{4\pi} \frac{I (d\mathbf{l} \times \mathbf{r})}{r^3}}
\end{equation}
where $\mu_0 = 4\pi \times 10^{-7}\,\mathrm{T}\cdot\mathrm{m}\cdot\mathrm{A}^{-1}$ is vacuum permeability.
\textbf{2. Magnetic Lorentz Force on a Moving Charge:}
\begin{equation}
  \mathbf{\mathbf{F} = q(\mathbf{E} + \mathbf{v} \times \mathbf{B})}
\end{equation}
\begin{itemize}
  \item Pure magnetic force: $\mathbf{F}_B = q(\mathbf{v} \times \mathbf{B})$. Since $\mathbf{F}_B \perp \mathbf{v}$, the instantaneous power is $P = \mathbf{F}_B \cdot \mathbf{v} = 0$.
  \item \textbf{Work Done by Magnetic Force is Identically Zero:} Magnetic forces change the direction of motion, never particle kinetic energy or speed.
  \item \textbf{Helical Trajectory in Uniform Field $\mathbf{B}$:}
  When particle velocity $\mathbf{v}$ makes angle $\theta$ with $\mathbf{B}$ ($v_\perp = v\sin\theta, v_\parallel = v\cos\theta$):
  \begin{equation}
    \mathbf{r = \frac{m v_\perp}{q B} = \frac{m v \sin\theta}{q B}}, \qquad \mathbf{T = \frac{2\pi m}{q B}}, \qquad \mathbf{\text{Pitch } p = v_\parallel T = \frac{2\pi m v \cos\theta}{q B}}
  \end{equation}
\end{itemize}
\end{conceptbox}

\begin{formulabox}[Master Magnetic Field Formulas: Straight Wires, Loops, Solenoids \& Toroids]
\textbf{1. Finite Straight Wire Subtending Angles $\theta_1, \theta_2$ at Distance $d$:}
\begin{equation}
  \mathbf{B = \frac{\mu_0 I}{4\pi d}(\sin\theta_1 + \sin\theta_2)}
\end{equation}
\begin{itemize}
  \item \textbf{Infinitely Long Straight Wire ($\theta_1 = \theta_2 = 90^\circ$):} $\mathbf{B = \frac{\mu_0 I}{2\pi d}}$.
  \item \textbf{Semi-Infinite Wire End ($\theta_1 = 90^\circ, \theta_2 = 0^\circ$):} $\mathbf{B = \frac{\mu_0 I}{4\pi d}}$.
\end{itemize}
\textbf{2. Circular Current Loop of Radius $R$ on Symmetry Axis ($x$):}
\begin{equation}
  \mathbf{B(x) = \frac{\mu_0 I R^2}{2(R^2 + x^2)^{3/2}} \hat{\mathbf{i}}} \implies \mathbf{B_{\text{center}} = \frac{\mu_0 I}{2R}}
\end{equation}
\textbf{3. Ampere's Circuital Law:}
\begin{equation}
  \mathbf{\oint \mathbf{B}\cdot d\mathbf{l} = \mu_0 I_{\text{enclosed}}}
\end{equation}
\begin{itemize}
  \item \textbf{Ideal Long Solenoid ($n$ turns/meter):} Interior field is axial and uniform: $\mathbf{B = \mu_0 n I}$. Near ends: $B_{\text{end}} = \frac{1}{2}\mu_0 n I$.
  \item \textbf{Toroid of Mean Radius $R$ and $N$ Total Turns:} $n = \frac{N}{2\pi R} \implies \mathbf{B = \mu_0 n I = \frac{\mu_0 N I}{2\pi R}}$.
\end{itemize}
\textbf{4. Force Between Parallel Currents (Definition of Ampere):}
\begin{equation}
  \mathbf{\frac{F}{L} = \frac{\mu_0 I_1 I_2}{2\pi d}}
\end{equation}
Parallel currents \textbf{attract}; antiparallel currents \textbf{repel}.
\end{formulabox}

\begin{formulabox}[Moving Coil Galvanometer Conversion \& Magnetic Materials]
\textbf{1. Moving Coil Galvanometer (MCG):}
Deflecting torque balances radial restoring suspension torque: $\tau = N I A B = C \theta \implies \theta = \left(\frac{NAB}{C}\right)I$.
\begin{itemize}
  \item \textbf{Conversion to Ammeter (Extending range to $I > I_g$):} Connect small shunt resistance $S$ in parallel:
  \begin{equation}
    (I - I_g) S = I_g R_g \implies \mathbf{S = \frac{I_g R_g}{I - I_g}}
  \end{equation}
  \item \textbf{Conversion to Voltmeter (Extending range to $V$):} Connect high multiplier resistance $R$ in series:
  \begin{equation}
    V = I_g (R_g + R) \implies \mathbf{R = \frac{V}{I_g} - R_g}
  \end{equation}
\end{itemize}
\textbf{2. Classification of Magnetic Materials:}
\begin{itemize}
  \item \textbf{Diamagnetic:} $\chi_m < 0$ (small, negative, independent of temperature). $\mu_r < 1$. Repelled by magnetic fields. Examples: $\mathrm{Bi, Cu, H_2O, NaCl}$.
  \item \textbf{Paramagnetic:} $\chi_m > 0$ (small, positive). Follows \textbf{Curie's Law}: $\chi_m = \frac{C}{T}$. $\mu_r > 1$. Weakly attracted. Examples: $\mathrm{Al, Na, O_2}$.
  \item \textbf{Ferromagnetic:} $\chi_m \gg 1$ (enormous, positive). Follows \textbf{Curie-Weiss Law} above Curie temperature $T_c$: $\chi_m = \frac{C}{T - T_c}$. Exhibits domains and hysteresis. Examples: $\mathrm{Fe, Co, Ni}$.
\end{itemize}
\end{formulabox}

\begin{trapbox}[Syllabus Boundary Alert: Cyclotron \& Earth's Magnetism Deletion]
\textbf{NMC REVISION WARNING FOR NEET UG:}
Under the revised NEET UG syllabus issued by the National Medical Commission:
\begin{itemize}
  \item \textbf{STRICTLY DELETED:} The \textbf{Cyclotron} (resonating RF frequency $\nu_c = \frac{qB}{2\pi m}$, kinetic energy $K_{\max} = \frac{q^2 B^2 R^2}{2m}$, relativistic limits) has been \textbf{completely removed}.
  \item \textbf{STRICTLY DELETED:} \textbf{Earth's Magnetic Elements} (Magnetic Declination $\theta$, Angle of Dip / Inclination $\delta$, Horizontal Component $B_H$, relation $B = \sqrt{B_H^2 + B_V^2}$) have been \textbf{completely removed}.
\end{itemize}
\end{trapbox}

\section{Electromagnetic Induction (EMI) \& Inductive Circuits}

\begin{conceptbox}[Faraday's Law, Lenz's Law \& Motional EMF]
\textbf{1. Magnetic Flux ($\Phi_B$):}
\begin{equation}
  \Phi_B = \iint \mathbf{B}\cdot d\mathbf{A} = B A \cos\theta \quad [\mathrm{Weber}, \mathrm{Wb} = \mathrm{T}\cdot\mathrm{m}^2]
\end{equation}
\textbf{2. Faraday's Law of Electromagnetic Induction:}
Induced EMF in any closed circuit equals the negative time rate of change of magnetic flux:
\begin{equation}
  \mathbf{\mathcal{E} = -\frac{d\Phi_B}{dt} = -N \frac{d\Phi_B}{dt}}
\end{equation}
\textbf{3. Lenz's Law (Conservation of Energy):}
The polarity of the induced EMF is such that it produces a current whose magnetic field opposes the original flux change that created it.
\textbf{4. Motional EMF Derivations:}
\begin{itemize}
  \item \textbf{Straight Conducting Rod of Length $L$ Moving with Velocity $\mathbf{v} \perp \mathbf{B}$:}
  Lorentz force on free electrons $\mathbf{F} = -e(\mathbf{v} \times \mathbf{B})$ drives charge accumulation until internal electrostatic field cancels magnetic force: $eE = evB \implies \mathbf{\mathcal{E} = B v L}$.
  \item \textbf{Conducting Rod of Length $L$ Rotating at Angular Velocity $\omega$ in Uniform $\mathbf{B}$:}
  Integrating elemental motional EMF $d\mathcal{E} = B v(r) dr = B (\omega r) dr$:
  \begin{equation}
    \mathbf{\mathcal{E} = \int_0^L B \omega r dr = \frac{1}{2} B \omega L^2}
  \end{equation}
\end{itemize}
\end{conceptbox}

\begin{formulabox}[Self-Inductance, Mutual Inductance \& Magnetic Energy Density]
\textbf{1. Self-Inductance ($L$):}
Total flux linkage is proportional to circulating current: $N \Phi_B = L I$.
\begin{equation}
  \mathbf{\mathcal{E}_L = -L \frac{dI}{dt}} \quad [\mathrm{Henry}, \mathrm{H} = \mathrm{V}\cdot\mathrm{s}\cdot\mathrm{A}^{-1}]
\end{equation}
For a long air-core solenoid of length $l$, area $A$, and turn density $n$:
\begin{equation}
  \Phi = B A = (\mu_0 n I) A \implies N \Phi = (n l)(\mu_0 n I A) = \mu_0 n^2 A l I \implies \mathbf{L = \mu_0 n^2 A l = \mu_0 \frac{N^2}{l} A}
\end{equation}
\textbf{2. Mutual Inductance ($M$):}
$N_2 \Phi_2 = M_{21} I_1 \implies \mathcal{E}_2 = -M \frac{dI_1}{dt}$. For coaxial solenoids:
\begin{equation}
  \mathbf{M = \mu_0 n_1 n_2 A l = \frac{\mu_0 N_1 N_2 A}{l}}
\end{equation}
Coupling coefficient: $\mathbf{k = \frac{M}{\sqrt{L_1 L_2}}}$ ($0 \le k \le 1$).
\textbf{3. Energy Stored in an Inductor:}
\begin{equation}
  \mathbf{U = \int_0^I L i di = \frac{1}{2} L I^2}
\end{equation}
\textbf{4. Magnetic Energy Density in a Magnetic Field:}
\begin{equation}
  \mathbf{u_B = \frac{U}{\text{Volume}} = \frac{B^2}{2\mu_0}}
\end{equation}
\end{formulabox}

\begin{ncertbox}[NCERT Extraction: Eddy Currents \& Practical Technological Applications]
When bulk conductors are subjected to changing magnetic flux, circulating closed induced current loops (\textbf{Eddy Currents}) are established in planes perpendicular to the magnetic lines of force.
\begin{itemize}
  \item \textbf{Electromagnetic Damping:} The metallic non-magnetic frame of a moving coil galvanometer develops opposing eddy currents as it deflects, rapidly bringing the pointer to rest at its steady reading without oscillatory hunting (dead-beat galvanometer).
  \item \textbf{Induction Furnaces:} High-frequency alternating magnetic fields generate massive eddy currents within metallic charges, producing Joule heat sufficient to melt metals in vacuum.
  \item \textbf{Magnetic Braking in Trains:} Electromagnets activated above rails induce opposing eddy currents in the rotating metallic wheels/rails, producing smooth, wear-free deceleration.
  \item \textbf{Mitigation in Transformers:} To minimize eddy current heating losses in transformers and motors, solid iron cores are replaced with thin, individually varnished/insulated metallic laminations stacked parallel to magnetic flux lines.
\end{itemize}
\end{ncertbox}

\section{Alternating Currents (AC) \& Resonant Circuits}

\begin{conceptbox}[Alternating Current Fundamentals: Phasors \& RMS Values]
A sinusoidal alternating EMF and current:
\begin{equation}
  V(t) = V_0 \sin(\omega t), \qquad I(t) = I_0 \sin(\omega t - \phi)
\end{equation}
\textbf{1. Mean (Average) Value over Half-Cycle:}
\begin{equation}
  I_{\text{avg}} = \frac{1}{\pi/\omega} \int_0^{\pi/\omega} I_0 \sin(\omega t) dt = \mathbf{\frac{2}{\pi} I_0 \approx 0.637 I_0}
\end{equation}
Full cycle average is strictly zero: $\langle I \rangle_{\text{full}} = 0$.
\textbf{2. Root-Mean-Square (RMS / Virtual) Value:}
Equating Joule heating over a full cycle $H = \int_0^T I^2 R dt = I_{\text{rms}}^2 R T$:
\begin{equation}
  I_{\text{rms}} = \sqrt{\frac{1}{T} \int_0^T I_0^2 \sin^2(\omega t) dt} = \mathbf{\frac{I_0}{\sqrt{2}} \approx 0.707 I_0}, \qquad \mathbf{V_{\text{rms}} = \frac{V_0}{\sqrt{2}}}
\end{equation}
AC meters (hot-wire ammeters and voltmeters) always read \textbf{RMS values}.
\end{conceptbox}

\begin{formulabox}[Series LCR Circuit: Impedance, Resonance \& Quality Factor]
In a series LCR circuit with supply voltage $V(t) = V_0 \sin(\omega t)$:
\begin{equation}
  V_R = I R, \qquad V_L = I X_L = I (\omega L), \qquad V_C = I X_C = I \left(\frac{1}{\omega C}\right)
\end{equation}
Since $V_L$ leads current by $\pi/2$ and $V_C$ lags current by $\pi/2$:
\begin{equation}
  V_0 = \sqrt{V_R^2 + (V_L - V_C)^2} = I_0 \sqrt{R^2 + \left(\omega L - \frac{1}{\omega C}\right)^2}
\end{equation}
\begin{equation}
  \mathbf{Z = \sqrt{R^2 + (X_L - X_C)^2} = \sqrt{R^2 + \left(\omega L - \frac{1}{\omega C}\right)^2}}, \qquad \mathbf{\tan\phi = \frac{X_L - X_C}{R}}
\end{equation}
\textbf{1. Electrical Resonance Condition:}
When inductive reactance equals capacitive reactance:
\begin{equation}
  X_L = X_C \implies \omega_0 L = \frac{1}{\omega_0 C} \implies \mathbf{\omega_0 = \frac{1}{\sqrt{LC}}}, \quad \mathbf{\nu_0 = \frac{1}{2\pi\sqrt{LC}}}
\end{equation}
At resonance:
\begin{itemize}
  \item Impedance is minimum and purely resistive: $\mathbf{Z_{\min} = R}$.
  \item Current amplitude is maximum and in phase with voltage: $\mathbf{I_{0,\max} = \frac{V_0}{R}}$, $\phi = 0$.
  \item Power factor is unity: $\cos\phi = 1$.
\end{itemize}
\textbf{2. Sharpness of Resonance \& Quality Factor ($Q$):}
\begin{equation}
  \mathbf{Q = \frac{\omega_0 L}{R} = \frac{1}{\omega_0 C R} = \frac{1}{R}\sqrt{\frac{L}{C}} = \frac{\omega_0}{2\Delta\omega} = \frac{\text{Resonant Frequency}}{\text{Bandwidth}}}
\end{equation}
High $Q$-factor indicates sharp selectivity and low damping.
\textbf{3. Average Power in AC Circuit \& Power Factor:}
\begin{equation}
  \mathbf{P_{\text{avg}} = V_{\text{rms}} I_{\text{rms}} \cos\phi}
\end{equation}
where $\cos\phi = \frac{R}{Z}$ is the \textbf{power factor}.
\begin{itemize}
  \item Pure $L$ or pure $C$: $\phi = \pm \pi/2 \implies \cos\phi = 0 \implies P_{\text{avg}} = 0$ (\textbf{Wattless Current} $I_w = I_{\text{rms}} \sin\phi$).
  \item Pure $R$ or LCR at resonance: $\phi = 0 \implies \cos\phi = 1 \implies P_{\text{avg}} = V_{\text{rms}} I_{\text{rms}}$.
\end{itemize}
\textbf{4. Transformers:}
\begin{equation}
  \mathbf{\frac{V_s}{V_p} = \frac{N_s}{N_p} = \frac{I_p}{I_s} = K} \quad (\text{Transformation Ratio})
\end{equation}
Efficiency: $\eta = \frac{P_{\text{output}}}{P_{\text{input}}} = \frac{V_s I_s \cos\phi_s}{V_p I_p \cos\phi_p}$.
\end{formulabox}

\section{Electromagnetic Waves \& Maxwell's Equations}

\begin{conceptbox}[Maxwell's Displacement Current \& Unified Field Equations]
\textbf{1. Maxwell's Displacement Current ($I_d$):}
Ampere's Circuital Law ($\oint \mathbf{B}\cdot d\mathbf{l} = \mu_0 I$) failed when applied to a charging capacitor boundary. Maxwell resolved this inconsistency by postulating that a time-varying electric field produces an equivalent displacement current:
\begin{equation}
  \Phi_E = E A = \left(\frac{Q}{\varepsilon_0 A}\right) A = \frac{Q}{\varepsilon_0} \implies I_d = \varepsilon_0 \frac{d\Phi_E}{dt} = \frac{dQ}{dt} = I_c
\end{equation}
\begin{equation}
  \mathbf{\oint \mathbf{B}\cdot d\mathbf{l} = \mu_0 (I_c + I_d) = \mu_0 I_c + \mu_0 \varepsilon_0 \frac{d\Phi_E}{dt}} \quad (\text{Ampere-Maxwell Law})
\end{equation}
\textbf{2. Transverse Nature \& Speed of EM Waves:}
Electric and magnetic vectors oscillate perpendicular to each other and perpendicular to the wave propagation direction:
\begin{equation}
  \mathbf{c = \frac{1}{\sqrt{\mu_0 \varepsilon_0}} = \frac{E_0}{B_0} \approx 3.00 \times 10^8\,\mathrm{m/s}}
\end{equation}
\textbf{3. Poynting Vector \& Radiation Pressure:}
Energy flux vector: $\mathbf{S = \frac{1}{\mu_0}(\mathbf{E} \times \mathbf{B})}$. Average intensity: $I = \frac{1}{2}\varepsilon_0 c E_0^2 = \frac{E_0 B_0}{2\mu_0}$.
\begin{equation}
  \text{Radiation Pressure: } P_{\text{rad}} = \begin{cases} \frac{I}{c} & \text{Perfectly Absorbing Surface} \\ \frac{2I}{c} & \text{Perfectly Reflecting Surface} \end{cases}
\end{equation}
\end{conceptbox}

\begin{formulabox}[The Electromagnetic Spectrum: Master Classification]
\begin{center}
\begin{tabularx}{\linewidth}{|l|c|c|Y|}
\hline
\textbf{Wave Type} & \textbf{Wavelength ($\lambda$)} & \textbf{Frequency ($\nu$)} & \textbf{Production Mechanism \& Major Use} \\
\hline
\textbf{Radio Waves} & $> 0.1\,\mathrm{m}$ & $< 3 \times 10^9\,\mathrm{Hz}$ & Rapid charge acceleration in LC antennas; Radio/TV communication \\
\hline
\textbf{Microwaves} & $1\,\mathrm{mm} - 0.1\,\mathrm{m}$ & $10^9 - 10^{11}\,\mathrm{Hz}$ & Klystron/Magnetron tubes; Radar navigation, Microwave ovens \\
\hline
\textbf{Infrared} & $700\,\mathrm{nm} - 1\,\mathrm{mm}$ & $10^{11} - 4 \times 10^{14}\,\mathrm{Hz}$ & Thermal molecular vibrations; Greenhouse effect, Remote controls \\
\hline
\textbf{Visible Light} & $400\,\mathrm{nm} - 700\,\mathrm{nm}$ & $4 - 7.5 \times 10^{14}\,\mathrm{Hz}$ & Electron transitions in outer atomic shells; Human vision \\
\hline
\textbf{Ultraviolet} & $1\,\mathrm{nm} - 400\,\mathrm{nm}$ & $10^{15} - 10^{17}\,\mathrm{Hz}$ & High-temperature inner shell transitions, sun; Sterilization, LASIK \\
\hline
\textbf{X-Rays} & $0.01\,\mathrm{nm} - 10\,\mathrm{nm}$ & $10^{16} - 10^{20}\,\mathrm{Hz}$ & Fast electrons decelerated on heavy metal target; Medical radiography \\
\hline
\textbf{Gamma Rays} & $< 10^{-3}\,\mathrm{nm}$ & $> 10^{19}\,\mathrm{Hz}$ & Nuclear de-excitation / nuclear isomeric transitions; Cancer radiotherapy \\
\hline
\end{tabularx}
\end{center}
\end{formulabox}

\section{Master High-Yield Problem Archetypes \& Worked Solutions}

\begin{usecasebox}[Archetype 1: Cube Resistor Network Equivalent Resistances]
\textbf{Problem:} Twelve identical resistors, each of resistance $R$, form the edges of a skeleton cube. Compute equivalent resistance between: (i) diagonally opposite body corners, (ii) opposite face corners, and (iii) adjacent edge corners.
\textbf{Step-by-Step Solution:}
\begin{enumerate}
  \item \textbf{Body Diagonal (Points $A$ and $G$):}
  Inject current $I$ at $A$ and extract at $G$. By threefold rotational symmetry around the body diagonal:
  \begin{itemize}
    \item 3 edges leave $A$, each carrying $I/3$.
    \item Each splits into 2 edges, each carrying $I/6$ across the middle 6 edges.
    \item These recombine into 3 edges entering $G$, each carrying $I/3$.
  \end{itemize}
  Total potential drop: $V_{AG} = \left(\frac{I}{3}\right)R + \left(\frac{I}{6}\right)R + \left(\frac{I}{3}\right)R = I R \left(\frac{1}{3} + \frac{1}{6} + \frac{1}{3}\right) = \frac{5}{6} I R$.
  \begin{equation}
    \mathbf{R_{\text{body}} = \frac{5}{6} R}
  \end{equation}
  \item \textbf{Face Diagonal (Points $A$ and $C$):}
  By mirror symmetry across the plane bisecting the face:
  \begin{equation}
    \mathbf{R_{\text{face}} = \frac{3}{4} R}
  \end{equation}
  \item \textbf{Adjacent Edge (Points $A$ and $B$):}
  \begin{equation}
    \mathbf{R_{\text{edge}} = \frac{7}{12} R}
  \end{equation}
\end{enumerate}
\end{usecasebox}

\begin{usecasebox}[Archetype 2: Series LCR Circuit Resonance \& Power Factor Calculation]
\textbf{Problem:} A series LCR circuit has $R = 10\,\Omega$, $L = 0.1\,\mathrm{H}$, and $C = 100\,\mu\mathrm{F}$ connected to an AC source $V(t) = 200 \sin(100 t)\,\mathrm{V}$. Calculate: (i) Inductive reactance $X_L$, (ii) Capacitive reactance $X_C$, (iii) Total impedance $Z$, (iv) Peak current $I_0$, (v) Resonant frequency $\omega_0$, and (vi) Quality factor $Q$.
\textbf{Step-by-Step Solution:}
\begin{enumerate}
  \item $\omega = 100\,\mathrm{rad/s}$, $V_0 = 200\,\mathrm{V}$.
  \item $X_L = \omega L = 100 \times 0.1 = \mathbf{10\,\Omega}$.
  \item $X_C = \frac{1}{\omega C} = \frac{1}{100 \times 100 \times 10^{-6}} = \frac{1}{10^{-2}} = \mathbf{100\,\Omega}$.
  \item Total impedance $Z$:
  \begin{equation}
    Z = \sqrt{R^2 + (X_L - X_C)^2} = \sqrt{10^2 + (10 - 100)^2} = \sqrt{100 + (-90)^2} = \sqrt{100 + 8100} = \sqrt{8200} \approx \mathbf{90.55\,\Omega}
  \end{equation}
  Circuit is strongly capacitive ($X_C > X_L$), current leads voltage.
  \item Peak current: $I_0 = \frac{V_0}{Z} = \frac{200}{90.55} \approx \mathbf{2.21\,\mathrm{A}}$.
  \item Resonant angular frequency $\omega_0$:
  \begin{equation}
    \omega_0 = \frac{1}{\sqrt{LC}} = \frac{1}{\sqrt{0.1 \times 10^{-4}}} = \frac{1}{\sqrt{10^{-5}}} = \frac{1000}{\sqrt{10}} \approx \mathbf{316.2\,\mathrm{rad/s}}
  \end{equation}
  \item Quality factor $Q$:
  \begin{equation}
    Q = \frac{1}{R}\sqrt{\frac{L}{C}} = \frac{1}{10}\sqrt{\frac{0.1}{10^{-4}}} = \frac{1}{10}\sqrt{1000} = \frac{31.62}{10} \approx \mathbf{3.16}
  \end{equation}
\end{enumerate}
\end{usecasebox}

\end{document}
